% !TEX program = xelatex
\documentclass[11pt]{article}

\usepackage{my-template}
\usepackage{microtype}

\definecolor{qaaccent}{HTML}{315EFB}
\definecolor{qasoft}{HTML}{EEF3FF}
\definecolor{qamuted}{HTML}{5F6773}

\hypersetup{
  colorlinks=true,
  linkcolor=qaaccent,
  urlcolor=qaaccent
}

\setlength{\parindent}{0pt}
\setlength{\parskip}{5pt}
\setlength{\tabcolsep}{4pt}
\renewcommand{\arraystretch}{1.18}
\renewcommand{\contentsname}{Mục lục}
\renewcommand{\tablename}{Bảng}

\newcommand{\quickanswer}[1]{%
  \par\medskip
  \noindent\colorbox{qasoft}{%
    \parbox{\dimexpr\linewidth-2\fboxsep\relax}{%
      \textcolor{qaaccent}{\textbf{Trả lời ngắn.}} #1}}
  \par\medskip
}

\begin{document}
\begin{center}
  {\LARGE\bfseries Q\&A kỹ thuật\par}
  \vspace{3pt}
  {\large Danh mục tra cứu nhanh cho nhóm\par}
  \vspace{5pt}
  {\small\textcolor{qamuted}{Cập nhật: 05/09/2026}\par}
  \vspace{6pt}
  \color{qaaccent}\rule{0.72\linewidth}{1.2pt}
\end{center}

\tableofcontents

\section{Giải mã Fairseq}

\subsection{Tại sao độ dài tối đa ở E1 và E3-custom được đặt bằng
\texorpdfstring{$1{,}2$}{1,2} lần độ dài đầu vào cộng 10 token?}

\quickanswer{Đây là một giới hạn kỹ thuật được chọn theo thống kê của corpus,
không phải một hằng số do bài báo về Transformer hay beam search đề xuất. Hệ số
$1{,}2$ bao phủ độ dài bản dịch hợp lệ trên các tập gold, còn phần cộng 10 bảo
vệ các câu nguồn ngắn và ngăn mô hình sinh lặp quá dài.}

\paragraph{Một token ở nhánh Fairseq là gì?}
Nhánh Fairseq không áp dụng BPE hay SentencePiece: không có tùy chọn
\texttt{--bpe} và cũng không có bước phân đoạn dưới từ trong
\texttt{fairseq\_runner.py} hoặc các cấu hình. Vì vậy, một ``token'' mà
decoder nhìn thấy chính là một đơn vị phân tách bằng khoảng trắng: phía nguồn
là âm tiết tiếng Việt, còn phía đích là từng Hán tự. Thống kê độ dài của các
tệp corpus do đó cũng chính là thống kê độ dài đầu vào và đầu ra của mô hình.

\begin{table}[H]
  \centering
  \scriptsize
  \caption{Thống kê độ dài thực tế của corpus.}
  \label{tab:qa-fairseq-length}
  \begin{tabularx}{\linewidth}{@{}lccccc>{\raggedright\arraybackslash}X@{}}
    \toprule
    Phân hoạch
      & \shortstack{Token nguồn\\TB / max}
      & \shortstack{Token đích\\TB / max}
      & \shortstack{$p_{50}$\\đích/nguồn}
      & $p_{95}$
      & Max
      & Tham chiếu vượt $1{,}2|x|+10$ \\
    \midrule
    Train (19.218)
      & 20,9 / 199
      & 18,6 / 196
      & 0,870
      & 1,417
      & 49,3
      & 353 (1,84\%) \\
    Test gold (510)
      & 20,7 / 90
      & 17,6 / 79
      & 0,864
      & 1,030
      & 1,190
      & 0 (0,00\%) \\
    \bottomrule
  \end{tabularx}
\end{table}

\paragraph{Vì sao chọn $a=1{,}2$?}
Hán văn cổ thường cô đọng hơn tiếng Việt. Trung vị độ dài bản dịch gold chỉ
xấp xỉ 0,86 token đích trên mỗi token nguồn; trên tập test gold, tham chiếu có
tỷ lệ lớn nhất là 1,190. Vì vậy, $1{,}2$ nằm ngay trên cực đại quan sát được và
cao hơn trung vị khoảng 40\%. Ngưỡng này đủ rộng để không cắt một bản dịch hợp
lệ trên test, nhưng đủ chặt để giới hạn vòng lặp sinh lặp. Nếu dùng giới hạn
mặc định dạng $a=0,b=200$, một đầu vào chỉ bốn token vẫn có thể sinh tới 200
Hán tự khi decoder không phát ra token kết thúc.

\paragraph{Vì sao cộng thêm $b=10$?}
Tỷ lệ độ dài biến động mạnh ở các câu ngắn. Với đầu vào hai token,
$1{,}2\times2=2{,}4$ sẽ cắt một bản dịch hoàn toàn bình thường dài bốn hoặc năm
Hán tự. Phần cộng 10 đóng vai trò khoảng đệm cho những câu ngắn, vốn xuất hiện
khá nhiều trong corpus, chẳng hạn ``趙 紀''. Ở các câu dài, thành phần
$1{,}2|x|$ chi phối nên khoảng đệm này hầu như không làm giới hạn trở nên quá
rộng.

\paragraph{Lợi ích thực tế.}

\begin{itemize}
  \item \textbf{Giảm thời gian giải mã.} Vòng lặp beam có thể chạy tới
  \texttt{max\_len} nếu một giả thuyết chưa sinh token kết thúc. Với beam 7 và
  batch 64, giới hạn khoảng 35 bước cho một câu điển hình nhanh hơn đáng kể so
  với cho phép tối đa 200 bước trên 510 câu và nhiều thí nghiệm.

  \item \textbf{Nhất quán giữa chọn checkpoint và giải mã cuối.} Bốn giá trị
  beam, length penalty, $a$ và $b$ được truyền vào
  \texttt{--eval-bleu-args} khi huấn luyện
  (\texttt{fairseq\_runner.py:166}) và lệnh generate
  (\texttt{fairseq\_runner.py:242}). Vì vậy, checkpoint tốt nhất theo BLEU
  validation được chọn bằng đúng chế độ giải mã dùng trên test.

  \item \textbf{Bổ sung cho giới hạn toàn cục.} Cấu hình
  \texttt{max\_target\_positions: 256} khiến Fairseq luôn chặn độ dài ở 255
  token. Quy tắc $1{,}2|x|+10$ là giới hạn theo từng câu, chặt hơn nhưng vẫn
  nằm dưới hàng rào toàn cục này.

  \item \textbf{Kiểm soát độ dài có chủ đích.} Giá trị
  \texttt{lenpen=1.0} không tăng hoặc giảm thêm mức phạt độ dài ngoài phép
  chuẩn hóa thông thường. Do đó, quy tắc affine trên là giới hạn cứng chính và
  cần được đặt rõ ràng thay vì kế thừa mức 200 token.
\end{itemize}

\paragraph{So sánh với nhánh QLoRA.}
E2 và E4 dùng
\[
  \texttt{max\_new\_tokens}
  =\min\!\left(512,\left\lceil2L_{\mathrm{src}}\right\rceil+32\right).
\]
Đây vẫn là một giới hạn affine theo độ dài nguồn, nhưng được nới rộng vì
tokenizer của LLM dùng subword và có thể cần nhiều token hơn để biểu diễn một
Hán tự. Hai backend vì vậy tuân theo cùng một nguyên tắc thiết kế, dù không thể
dùng chung các hệ số.

\paragraph{Lưu ý về tập train silver.}
Trong tập train, 353 tham chiếu (1,84\%) vượt ngưỡng $1{,}2|x|+10$; tỷ lệ cực
đại 49,3 đến từ một số ít cặp tự căn chỉnh sai lệch. Điều này không trực tiếp
cắt dữ liệu huấn luyện vì giới hạn chỉ được áp dụng khi sinh câu. Trên các tập
gold, không có tham chiếu nào vượt ngưỡng, nên các ngoại lệ ở train phản ánh
nhiễu căn chỉnh hơn là bằng chứng cho thấy giới hạn giải mã quá chặt.

\section{Token và biểu diễn văn bản}

\subsection{Token trong hệ thống này là gì?}

\quickanswer{Từ \emph{token} mang ba nghĩa khác nhau tùy giai đoạn: đơn vị phân
tách bằng khoảng trắng trong mô hình Fairseq, subword của tokenizer Qwen3
trong mô hình LLM, và từng Hán tự trong bước chấm điểm. Không nên dùng lẫn ba
khái niệm này khi diễn giải giới hạn độ dài hoặc điểm BLEU.}

\begin{table}[H]
  \centering
  \scriptsize
  \caption{Đơn vị token tại ba giai đoạn của pipeline.}
  \label{tab:qa-token-stages}
  \begin{tabularx}{\linewidth}{
    @{}
    >{\raggedright\arraybackslash}p{2.45cm}
    >{\raggedright\arraybackslash}p{3.05cm}
    >{\raggedright\arraybackslash}p{1.75cm}
    >{\raggedright\arraybackslash}p{1.9cm}
    >{\raggedright\arraybackslash}X
    @{}}
    \toprule
    Giai đoạn & Đơn vị & Tiếng Việt & Hán văn & Chi phối \\
    \midrule
    Mô hình Fairseq (E1/E3-custom)
      & Đơn vị phân tách bằng khoảng trắng
      & Âm tiết
      & Một Hán tự
      & \texttt{max\_len\_a/b}; \texttt{max\_source\_positions: 256};
        \texttt{max\_tokens\_per\_gpu: 1000} \\
    Mô hình LLM (E2/E4)
      & Subword BPE của Qwen3
      & Không cố định
      & Không cố định
      & $\texttt{max\_new\_tokens}
        =\min(512,\lceil2L_{\mathrm{src}}\rceil+32)$ \\
    Chấm điểm
      & Một ký tự Unicode
      & Không áp dụng
      & Một Hán tự
      & \texttt{score\_form\_chinese} và SacreBLEU 13a \\
    \bottomrule
  \end{tabularx}
\end{table}

\paragraph{Token của Fairseq.}
Không có BPE trong nhánh E1/E3-custom, nên Fairseq đọc trực tiếp các đơn vị
được ngăn cách bằng khoảng trắng trong corpus. Kiểm tra trên dữ liệu cho thấy
357.806 token trong các tệp \texttt{.cn} đều dài đúng một ký tự, thuộc 5.386
loại ký tự. Các tệp \texttt{.vi} đã được chuyển về chữ thường, bỏ dấu câu và
tách theo âm tiết, gồm 401.997 token thuộc 3.607 loại. Chẳng hạn,
\emph{viên ngoại lang} được tính là ba token, không phải một từ ghép duy nhất.

Các con số này là đơn vị thực tế được dùng cho
\texttt{max\_source\_positions}, \texttt{max\_len\_a/b} và ngân sách
\texttt{max\_tokens\_per\_gpu}. Vì không có bước phân đoạn trung gian, thống
kê độ dài trực tiếp trên tệp \texttt{.vi}/\texttt{.cn} cũng chính là thống kê
mà encoder và decoder Fairseq nhìn thấy.

\paragraph{Vì sao phía Hán văn dùng từng ký tự?}
Hán văn cổ vốn được viết liền, không có khoảng trắng xác định ranh giới từ.
Các bộ tách từ có sẵn như Jieba hoặc PKUSeg chủ yếu được huấn luyện trên tiếng
Trung hiện đại; áp dụng trực tiếp cho văn ngôn có thể áp đặt các ranh giới từ
hiện đại không phù hợp. Trong nhiều trường hợp, một Hán tự đồng thời mang vai
trò gần với một hình vị, nên biểu diễn cấp ký tự là lựa chọn phù hợp với ngữ
liệu này, thay vì chỉ là phương án dự phòng.

Biểu diễn trên cũng giữ từ vựng đích ở quy mô nhỏ. Với ngưỡng tần suất bằng 1,
mọi Hán tự xuất hiện trong tập train đều được giữ lại, nên không có
\texttt{<unk>} do cắt bỏ token hiếm trong chính tập train. Tuy nhiên, ký tự chỉ
xuất hiện ở validation hoặc test vẫn có thể nằm ngoài từ điển Fairseq; việc
không dùng BPE không loại bỏ dạng OOV này.

\paragraph{Vì sao phía tiếng Việt dùng âm tiết?}
Chính tả tiếng Việt đã đặt khoảng trắng giữa các âm tiết, nên thao tác
\texttt{split()} tạo ra một tokenizer xác định, không cần mô hình phụ hoặc
dependency bên ngoài. Tách từ đầy đủ sẽ cần thêm công cụ như VnCoreNLP hoặc
Underthesea, kéo theo một phiên bản phần mềm phải cố định và thêm nguồn sai số
tiền xử lý.

Trong corpus này, tiếng Việt mang nhiều âm đọc Hán--Việt, nên từng âm tiết
thường ánh xạ khá trực tiếp sang một Hán tự. Giữ đơn vị âm tiết bảo toàn cấu
trúc căn chỉnh gần đơn điệu và gần một--một, qua đó đơn giản hóa bài toán cho
Transformer sáu lớp được huấn luyện trên khoảng 19 nghìn câu. Quan hệ này cũng
giải thích vì sao trung vị tỷ lệ độ dài đích/nguồn xấp xỉ 0,86 thay vì biến
động khó dự đoán.

\paragraph{Vì sao Fairseq không dùng BPE?}
BPE thường được dùng để giới hạn kích thước từ vựng và xử lý từ chưa gặp. Ở
đây, từ vựng train của hai phía đã nhỏ, khoảng 3,6 nghìn loại tiếng Việt và
5,4 nghìn Hán tự. Các phép gộp subword học từ chỉ 19 nghìn câu có thể tạo thêm
một artifact tiền xử lý phải lưu phiên bản, đồng thời làm mờ ánh xạ gần
một--một giữa âm tiết Hán--Việt và Hán tự. Cấu hình
\texttt{threshold\_source: 1} và \texttt{threshold\_target: 1} giữ mọi loại
token có mặt trong train, nên không có token nào bị đổi thành \texttt{<unk>}
chỉ vì tần suất thấp.

\paragraph{Token của Qwen3.}
E2/E4 không đọc các đơn vị khoảng trắng nói trên như token mô hình. Toàn bộ
prompt và câu trả lời được tokenizer BPE của Qwen3 phân đoạn thành subword;
một âm tiết tiếng Việt hoặc một Hán tự có thể tương ứng với số token khác
nhau. Vì vậy, số token và giới hạn sinh của nhánh LLM không thể so sánh trực
tiếp với E1/E3-custom. Đây là lý do nhánh này dùng giới hạn rộng hơn:
\[
  \texttt{max\_new\_tokens}
  =\min\!\left(512,\left\lceil2L_{\mathrm{src}}\right\rceil+32\right).
\]

\paragraph{Token dùng để chấm điểm.}
Bước đánh giá cố ý định nghĩa một đơn vị thứ ba, độc lập với tokenizer của
từng mô hình. Hàm \texttt{score\_form\_chinese}
(\texttt{normalize.py:20}) xóa toàn bộ khoảng trắng khỏi dự đoán, sau đó chèn
lại đúng một khoảng trắng giữa từng ký tự. SacreBLEU với tokenizer 13a vì thế
tính các n-gram trên chuỗi Hán tự đã được chuẩn hóa; chrF++ cũng được tính với
\texttt{whitespace=False} trên cùng nội dung.

Việc chuẩn hóa này làm cho đầu ra Fairseq vốn có khoảng trắng và đầu ra Qwen3
vốn viết liền trở thành cùng một biểu diễn trước khi chấm. Tại
\texttt{evaluation.py:36--38}, pipeline còn tự tạo lại
\texttt{prediction\_scored} và từ chối tệp nếu kết quả không khớp với trường
đã lưu. Nhờ đó, một hệ thống không thể vô tình được chấm bằng cách token hóa
khác các hệ thống còn lại.

\paragraph{Ý nghĩa đối với phép so sánh E1 và E2.}
Hai backend không dùng chung token ở cấp mô hình, nên phép so sánh chỉ có ý
nghĩa khi cả hai được đưa về một token chấm điểm chung mà không backend nào
được ưu tiên. Trong thí nghiệm này, đơn vị chung đó là từng Hán tự. Điều này
cũng nối lại câu hỏi trước: vì Fairseq không dùng BPE, các thống kê trên tệp
thô đúng bằng đơn vị mà decoder Fairseq đếm, nên có thể kiểm tra trực tiếp
ngưỡng $1{,}2|x|+10$ trên corpus.

\section{Lựa chọn checkpoint}

\subsection{Tại sao E1 và E3-custom chọn theo BLEU validation, còn E2 chọn
theo validation loss?}

\quickanswer{E1 và E3-custom chọn checkpoint bằng BLEU validation với đúng cấu
hình beam search dùng khi đánh giá cuối, nên tiêu chí lựa chọn gắn trực tiếp
với chất lượng chuỗi dịch. E2 dùng response-only validation loss theo quy trình
chuẩn của Hugging Face Trainer vì phép tính teacher forcing đơn giản và rẻ hơn
sinh toàn bộ bản dịch. Tuy nhiên, lý do này là diễn giải hồi cứu; không có tài
liệu cùng thời điểm chạy hoặc ablation xác nhận quyết định đó.}

\begin{table}[H]
  \centering
  \small
  \caption{Khác biệt trong tiêu chí lựa chọn checkpoint.}
  \label{tab:qa-checkpoint-selection}
  \begin{tabularx}{\linewidth}{
    @{}
    >{\raggedright\arraybackslash}p{2.4cm}
    >{\raggedright\arraybackslash}p{3.2cm}
    >{\raggedright\arraybackslash}p{3.2cm}
    >{\raggedright\arraybackslash}X
    @{}}
    \toprule
    Hệ thống & Tiêu chí & Cách đánh giá & Hạn chế chính \\
    \midrule
    E1, E3-custom
      & BLEU validation lớn nhất
      & Sinh tự hồi quy bằng cùng cấu hình beam search dùng ở test
      & Chi phí sinh chuỗi cao hơn tính loss \\
    E2
      & Response-only validation loss nhỏ nhất
      & Teacher forcing trong Hugging Face Trainer
      & Không trực tiếp đo chất lượng dịch khi mô hình tự sinh \\
    \bottomrule
  \end{tabularx}
\end{table}

\paragraph{E1 và E3-custom.}
Hai mô hình Fairseq được chọn theo BLEU trên tập validation. Mỗi checkpoint
được giải mã bằng cùng kích thước beam, length penalty và giới hạn độ dài sẽ
được dùng ở bước đánh giá cuối. Cách làm này đặt tiêu chí lựa chọn sát với mục
tiêu chính của hệ thống: chất lượng toàn bộ chuỗi dịch khi mô hình tự sinh,
thay vì chỉ dựa trên xác suất của từng token khi đã biết tiền tố đúng.

Việc sinh validation định kỳ vẫn khả thi về tính toán vì E1 và E3-custom là
các Transformer tương đối nhỏ. Lựa chọn bằng BLEU cũng tránh phụ thuộc vào
cross-entropy validation có label smoothing của Fairseq. Giá trị loss này hữu
ích để theo dõi tối ưu, nhưng thứ hạng checkpoint theo loss không nhất thiết
trùng với thứ hạng theo BLEU cấp chuỗi.

\paragraph{E2.}
E2 chọn checkpoint có response-only validation loss nhỏ nhất thông qua quy
trình chuẩn của Hugging Face Trainer. Loss chỉ được tính trên token của câu trả
lời; các token thuộc system prompt và câu nguồn đã bị mask. Phép đánh giá dùng
teacher forcing, nên toàn bộ câu đích đúng được cung cấp cho mô hình và loss có
thể được tính trong một lượt forward, không cần chạy beam search.

Nếu chọn E2 bằng BLEU, pipeline phải bổ sung một vòng đánh giá riêng: dựng lại
prompt đúng định dạng chat, sinh tự hồi quy cho toàn bộ tập validation, tách
phần trả lời, chuẩn hóa Hán tự rồi mới gọi bộ chấm điểm. Với mô hình nhân quả
tám tỷ tham số và beam search, quy trình này tốn thời gian hơn đáng kể so với
tính validation loss bằng teacher forcing.

Tuy nhiên, repository không lưu một giải thích cùng thời điểm chạy cho lựa
chọn này và cũng không có ablation so sánh checkpoint tốt nhất theo loss với
checkpoint tốt nhất theo BLEU. Vì vậy, thuận lợi về tính toán và khả năng tích
hợp với Trainer chỉ là cách tái dựng hợp lý từ implementation, không nên được
trình bày như một quyết định thiết kế đã được xác minh.

\paragraph{Hệ quả đối với phép so sánh.}
Tiêu chí lựa chọn checkpoint chưa được đồng nhất hoàn toàn giữa ba hệ thống.
Validation loss không label smoothing của E2 là một mục tiêu token-level hợp
lệ trên dữ liệu giữ riêng, nhưng không đo trực tiếp chất lượng khi mô hình tự
sinh và có thể xếp hạng checkpoint khác với BLEU. Ngược lại, BLEU phản ánh đầu
ra tự hồi quy nhưng phụ thuộc vào cấu hình giải mã và tốn thêm chi phí sinh.

Tất cả quyết định lựa chọn đều chỉ sử dụng validation; tập test không trực
tiếp tham gia chọn checkpoint. Do đó, sự khác biệt này không phải rò rỉ test,
nhưng vẫn là một yếu tố gây nhiễu khi so sánh E1/E3-custom với E2. Một thí
nghiệm được kiểm soát tốt hơn nên chấm BLEU validation cho mọi checkpoint E2
đã lưu, rồi báo cáo xem tiêu chí BLEU và validation loss có chọn cùng một
checkpoint hay không.

\section{Truy hồi tri thức trong E3-custom}

\subsection{Câu đầu vào của E3-custom được tạo như thế nào?}

\quickanswer{E3-custom không sao chép 64 token từ câu dịch tham chiếu. Trong
bước tiền xử lý, hệ thống truy hồi tối đa 64 Hán tự ứng viên từ
\texttt{knowledge.json}, xếp hạng chúng bằng thông tin câu nguồn và tần suất
chỉ tính trên tập train, rồi nối chúng vào cuối câu tiếng Việt sau token phân
cách \texttt{\_\_knowledge\_\_}.}

\paragraph{Bước 1: xây dựng chỉ mục tra cứu.}
Mỗi mục Hán tự trong \texttt{knowledge.json} được liên kết với các dạng tra
cứu tiếng Việt lấy từ bốn trường:
\begin{itemize}
  \item \texttt{pronunciation};
  \item \texttt{nom\_query};
  \item \texttt{lookup\_query};
  \item \texttt{query\_variants}.
\end{itemize}
Các dạng này được chuẩn hóa rồi dùng làm khóa ánh xạ tới một hoặc nhiều Hán
tự. Chẳng hạn, dạng tiếng Việt \emph{lê} có thể trỏ tới nhiều ứng viên, trong
đó có ký tự 黎.

\paragraph{Bước 2: tìm các đoạn khớp trong câu nguồn.}
Retriever trượt qua mọi đoạn token liên tiếp của câu tiếng Việt và kiểm tra
xem đoạn đó có trùng với một dạng tra cứu trong chỉ mục hay không. Một đoạn
tiếng Việt có thể đề xuất nhiều Hán tự; ngược lại, cùng một Hán tự có thể được
đề xuất bởi nhiều đoạn khác nhau. Vì vậy, khoảng 20 đoạn khớp trong một câu có
thể tạo ra hơn 60 ứng viên phân biệt.

\paragraph{Bước 3: tính điểm từng ứng viên.}
Sau khi gộp các ứng viên trùng nhau, điểm chính của một ký tự $c$ là số đoạn
nguồn đã truy hồi ra ký tự đó:
\begin{equation}
  \operatorname{score}(c)
  =\operatorname{match\_count}(c).
\end{equation}
Một ký tự được nhiều đoạn nguồn cùng đề xuất sẽ có độ ưu tiên cao hơn ký tự
chỉ xuất hiện từ một lần khớp.

\paragraph{Bước 4: xếp hạng xác định.}
Các ứng viên được sắp xếp lần lượt theo:
\[
  \text{số lần khớp giảm dần}
  \;\rightarrow\;
  \text{tần suất đích trong train giảm dần}
  \;\rightarrow\;
  \text{thứ tự Unicode}.
\]
Tần suất Hán tự trong câu đích train chỉ đóng vai trò phá hòa khi hai ứng viên
có cùng số lần khớp. Câu dịch tham chiếu của mẫu đang xử lý không được xem;
câu đích validation và test cũng không tham gia xây chỉ mục tần suất. Quy tắc
Unicode cuối cùng giúp kết quả hoàn toàn xác định khi hai tiêu chí đầu vẫn
bằng nhau.

\paragraph{Bước 5: nối ứng viên vào câu nguồn.}
Sau khi xếp hạng, hệ thống lấy các ứng viên đứng đầu và tạo chuỗi:
\[
  \underbrace{\text{câu tiếng Việt}}_{\text{nguồn gốc}}
  \quad
  \texttt{\_\_knowledge\_\_}
  \quad
  c_1\;c_2\;\cdots\;c_k,
\]
trong đó
\begin{equation}
  k=\min\!\left(
    64,\;
    \text{số ứng viên truy hồi},\;
    256-L_{\mathrm{src}}-1
  \right).
\end{equation}
Thành phần cuối dành một vị trí cho token phân cách và bảo đảm chuỗi mở rộng
không vượt giới hạn 256 vị trí của encoder E3-custom.

Ví dụ, một câu nguồn sau khi bổ sung tri thức có dạng:
\begin{quote}
  \small\ttfamily
  tôn mẹ là lê thị làm linh hiển\\
  \_\_knowledge\_\_\\
  黎 是 宗 令 侍 氏 尊 孫 母 視 靈 羅 \ldots
\end{quote}
Các Hán tự được phân tách bằng khoảng trắng, nên Fairseq coi mỗi ứng viên là
một token phía nguồn. Token \texttt{\_\_knowledge\_\_} đánh dấu ranh giới
giữa câu tiếng Việt gốc và khối gợi ý, cho phép encoder học cách sử dụng hoặc
bỏ qua phần tri thức bổ sung.

\paragraph{Số lượng ứng viên thực tế.}
Không phải câu nào cũng nhận đúng 64 ứng viên, vì số lượng còn phụ thuộc vào
kết quả truy hồi và số vị trí encoder còn trống.
\begin{table}[H]
  \centering
  \small
  \caption{Số Hán tự ứng viên được nối vào mỗi câu.}
  \label{tab:qa-e3-candidates}
  \begin{tabular}{lrr}
    \toprule
    Phân hoạch & Trung bình & Lớn nhất \\
    \midrule
    Train & 61,53 & 64 \\
    Validation & 62,41 & 64 \\
    Test & 61,98 & 64 \\
    \bottomrule
  \end{tabular}
\end{table}
Mọi câu đều nhận ít nhất một gợi ý và phần lớn gần chạm trần. Điều này cho
thấy retriever có độ phủ cao nhưng khả năng chọn lọc thấp: nhiều Hán tự được
nối vào có thể không liên quan tới bản dịch cần sinh. Ngoài ra, một số ký tự
gợi ý chỉ xuất hiện ở validation hoặc test sẽ trở thành \texttt{<unk>} nếu
chúng không có trong từ vựng nguồn được xây từ train.

\paragraph{Vì sao giới hạn là 64?}
Không có tài liệu ghi lại lý do lựa chọn hoặc ablation cho giá trị 64. Cấu hình
\texttt{configs/e3\_custom\_fairseq\_knowledge.yaml:12} chỉ cố định
\texttt{max\_candidates: 64}; việc khớp, xếp hạng và nối ứng viên được triển
khai tại \texttt{code/mt\_pipeline/knowledge.py:80}. Do đó, 64 nên được mô tả
là một ngân sách ứng viên cố định của lần chạy, không phải giá trị đã được tối
ưu thực nghiệm.

\end{document}